\section{Mutação}

Cobertura mede o que foi \emph{executado}, não o que foi \emph{verificado}:
um teste sem asserção dá 100\% de cobertura
(\citetitle{artigo:testcoverage}, \citetitle{artigo:assertion-free-testing}).

O PIT introduz pequenos defeitos no código (por exemplo, \texttt{<=} vira
\texttt{<}) e roda os testes. Mutante que sobrevive costuma ser uma borda
que ninguém especificou. Em Maven, o PIT roda com
\texttt{mvn test-compile org.pitest:pitest-maven:mutationCoverage} e o
escopo se limita com \texttt{targetClasses}
(\citetitle{doc:pitest-maven}).

Limites da técnica:
\begin{itemize}
\item mutantes equivalentes (a mudança não altera o comportamento, então
  nenhum teste a detecta);
\item custo de execução;
\item o escore virando meta (lei de Goodhart: quando a medida vira alvo,
  deixa de ser boa medida).
\end{itemize}

Meça o resultado, como propõe o DORA (taxa de falha de mudanças, tempo de
recuperação; \citetitle{artigo:dora}), e use o escore como diagnóstico.
